	\documentclass[10pt,twoside]{article}
\usepackage[a5paper,inner=1.5cm,outer=2.0cm,top=2.0cm,bottom=2.5cm,pdftex]{geometry}
\usepackage[utf8]{inputenc}						% UTF-8 encoding
\usepackage[T1]{fontenc}
\usepackage[magyar,american]{babel}
\usepackage{csquotes}
\usepackage[backend=biber,
			hyperref=true,
			autolang=hyphen,
			style=numeric-comp,
			isbn=false,
			sorting=none,
			defernumbers=true,
			maxnames=10,
			doi=false,
			url=false,
			useprefix=true,
			]{biblatex}
\usepackage{palatino}							% Font
\usepackage{amsmath, amssymb, amsthm} 							% Theorem customization
\usepackage{graphicx,subfigure,overpic}
\usepackage{rotating}
\usepackage{multirow}
\usepackage{xcolor}
\usepackage[bf,small]{caption}
\addbibresource{../LatexProject/dissertation.bib}
\addbibresource{booklet.bib}

\newcount\posveccount
\newcommand*\posvec[1]{
        \global\posveccount#1
        [
        \posvecnext
}
\def\posvecnext#1{
        #1
        \global\advance\posveccount-1
        \ifnum\posveccount>0
                ,\
                \expandafter\posvecnext
        \else
                ]^{\mathrm{T}}
        \fi
} 
\newcommand{\vxs}{\mathbf{x}_{\mathrm{s}}}
\newcommand{\vxref}{\mathbf{x}_{\mathrm{ref}}}
\newcommand{\vk}{\mathbf{k}} 

\usepackage[pdftitle={A hangtérszintézis elmélet általánosítása (Tézisfüzet)},%
  pdfauthor={Firtha Gergely},% 
%  colorlinks=true,% 
%  urlcolor=blue,% 
%  linkcolor=blue,% 
%  citecolor=blue,% 
%  filecolor=green,% 
%  raiselinks=false,% 
  hyperfootnotes=true]{hyperref}

% Citation categories
\DeclareBibliographyCategory{journal} 
\DeclareBibliographyCategory{conference}
\DeclareBibliographyCategory{tech}
\DeclareBibliographyCategory{other}
% Citation commands
\newcommand*{\citeJ}[2][]{\addtocategory{journal}{#2}\cite[#1]{#2}}
\newcommand*{\citeC}[2][]{\addtocategory{conference}{#2}\cite[#1]{#2}}
\newcommand*{\citeO}[2][]{\addtocategory{other}{#2}\cite[#1]{#2}}
\newcommand*{\citeT}[2][]{\addtocategory{tech}{#2}\cite[#1]{#2}}
% No citation commands
\newcommand*{\nociteJ}[1]{\addtocategory{journal}{#1}\nocite{#1}}
\newcommand*{\nociteC}[1]{\addtocategory{conference}{#1}\nocite{#1}}
\newcommand*{\nociteO}[1]{\addtocategory{other}{#1}\nocite{#1}}
\newcommand*{\nociteT}[1]{\addtocategory{tech}{#1}\nocite{#1}}

\nociteJ{firtha2015sound,firtha2016wave_booklet,Firtha2016_booklet,doi:10.1121/1.4996126_booklet,schultz2017wave_booklet,Firtha2018:WFS_vs_SDM_booklet}
\nociteC{Firtha2012:isma_booklet,Firtha2013:daga_booklet,Firtha2013:internoise_booklet,Firtha2014:daga_booklet,Firtha2014:isma_booklet,
Firtha2015:daga_booklet,firtha2016:daga_booklet,Firtha2017:daga_booklet,Firtha2018_daga_a_booklet,Firtha2018_daga_moving_source_booklet}

\renewcommand{\bibfont}{\normalfont\small}
\renewcommand*{\mkbibnamegiven}[1]{%
\ifitemannotation{highlight}
{\textbf{#1}}
{#1}}



\renewcommand*{\mkbibnamefamily}[1]{%
\ifitemannotation{highlight}
{\textbf{#1}}
{#1}}

\newtheoremstyle{thesisgroupstyle}%
	{1em}% 				% Space above
	{5mm}% 				% Space below
	{\slshape}%			% Font type of body
	{0mm}%				% Indentation			
	{\bfseries}%		% Font of title
	{\newline}%			% Punctuation between head and body
	{1em}%				% Space after theorem head
	{}% 				% Manually specify head

\makeatother
\theoremstyle{thesisgroupstyle}
\newtheorem{thesisgroup}{Téziscsoport}
\renewcommand{\thethesisgroup}{\Roman{thesisgroup}}

\makeatletter
\newtheoremstyle{indented}
  {5pt}% space before
  {5pt}% space after
  {\addtolength{\@totalleftmargin}{3.5em}
   \addtolength{\linewidth}{-3.5em}
   \parshape 1 3.5em \linewidth}% body font
  {}% indent
  {\bfseries}% header font
  {.}% punctuation
  {.75em}% after theorem header
  {}% header specification (empty for default)
\makeatother
\theoremstyle{indented}
\newtheorem{thesis}{tézis}[thesisgroup]


\newcommand{\theauthor}{Firtha Gergey}
\newcommand{\authorprof}{M.Sc.E.E.}
\author{\theauthor}
\title{A hangtérszintézis elméletének általánosítása\\[.5ex]és alkalmazása mozgó források szintézisére}

\assignrefcontextkeyws[labelprefix=J]{J}
\assignrefcontextkeyws[labelprefix=C]{C}
\assignrefcontextkeyws[labelprefix=O]{O}

% Add underline for my author name
\def \mylastname {Firtha}
\newbibmacro{name:underlined}{%
    \ifthenelse{\equal{\namepartfamily}{\mylastname}}%
{\underline{\ifblank{\namepartgiven}{}{\namepartgiven\space}\namepartfamily}}%
{\ifblank{\namepartfamily}{}{\namepartgiven\space}\namepartfamily}%
    \ifthenelse{\value{listcount}<\value{liststop}}%
    {\addcomma\space}%
    {}%
}

\DeclareNameFormat{author}{%
    \nameparts{#1}% split the name data, will not be necessary in future versions
    \usebibmacro{name:underlined}%
    \usebibmacro{name:andothers}%
}


\begin{document}
\sloppy
\newcommand{\schoolname}{Budapesti Műszaki és Gazdaságtudományi Egyetem}
\newcommand{\facultyname}{Villamosmérnöki és Informatikai Kar}
\newcommand{\doctoralname}{Villamosmérnöki Tudományok Doktori Iskola}
\newcommand{\bookname}{PhD tézisfüzet}
\newcommand{\authorname}{Author}
\newcommand{\supervisorname}{Konzulens:}
\newcommand{\supervisor}{Dr. Fiala Péter}
\newcommand{\creationdate}{Budapest, 2018.}
\def \deptlogo {logos/hit_logo_en.png}
\def \lablogo {logos/last_logo_eng.png}

\selectlanguage{magyar}
\nonfrenchspacing

%\thispagestyle{empty}
\newcommand{\HRule}{\rule{\linewidth}{1pt}}

\begin{titlepage}
\begin{center}

%%% Header with university logo
\includegraphics[height=1.2cm]{logos/BME1782.pdf}
\HRule\par
{\small \textsc{ \bfseries \schoolname \\
\facultyname \\
\doctoralname}}\par

\vfill

%%% Inserting the title
{
\Large
\makeatletter
\@title
\makeatother
} \par
\bigskip
\bookname \par
\bigskip
{
\normalsize
\theauthor\\[.5ex]
\authorprof
}\par

\vfill
%%% Inserting author, supervisor and co-promotor
\begin{tabular}{@{}p{3cm}l@{}}
	\supervisorname & \supervisor \\
\end{tabular}

\vspace{\fill}

\includegraphics[height=1.2cm]{\deptlogo}\par

\vspace*{\stretch{0.33}}
%%% Inserting date
\normalsize
\creationdate



\end{center}


\end{titlepage}

\thispagestyle{empty}
\cleardoublepage
\setcounter{page}{1}

\section{Bevezetés}

\subsection{A hangtér-reprodukció célja}
%
A térhangzású technikák alapvető célja egy virtuális akusztikai környezet előállítása úgy, hogy a térben elhelyezett audio objektumokat a hallgató a kívánt térbeli jellemzőkkel (pl. virtuális hangforrás pozíciója, szélessége) érzékeli.
A hangtér-reprodukció ezt egy fix elhelyezésű hangszórósokaság megfelelő vezérlésével éri el azáltal, hogy a hangszórók által keltett hullámok szuperpozíciója a hallgató számára a kívánt akusztikai környezet hatását kelti.
Ezzel szemben a binauralizációs módszerek a térhangzást fejhallgatók jelének megfelelő szűrésével érik el.

A mai hangtér-reprodukciós rendszerek kiindulópontja a Blumlein által 1931-ben bevezetett kéthangszórós rendszer volt, bevezetve ezzel a sztereofónia alapjait.
A sztereofónia manapság is a legelterjedtebb reprodukciós módszer, olyan hangrendszereknek képezi a működési elvét, mint a Dolby 5.1, 7.1, az NHK 22.2 rendszere, vagy épp a jelenleg state-of-the-art kereskedelemben is kapható térhatású rendszerek, a Dolby Atmosnak és a DTS-X.
Általánosan véve---függetlenül a hangrendszerben alkalmazott hangszórók számától---a sztereofónia célja pusztán a térérzetet befolyásoló tényezők (pl. füljelek közötti késleltetés, illetve amplitúdókülönbség) visszaállítása a hallgató pozíciójában.
Ebből kifolyólag azonban a tökéletes térhangzás csak a tér egyetlen pontjában és szűk környezetében biztosítható, ez az ún. \emph{sweet spot}.
A fenti megszorítás a sztereofon technikák fontos közös hátrányának tekinthető.

A sztereofóniával ellentétben a fizikai hangtér-reprodukció (sound field synthesis) célja egy tetszőleges célhangtér fizikai jellemzőinek visszaállítása egy kiterjedt megfigyelési területen.
Amennyiben ezt sikerül elérni, a hallgató a megfigyelési terület tetszőleges pontján eleve a kívánt akusztikai érzeti jellemzőket érzékeli.
Természetesen egy hullámtér kiterjedt területen való szabályozása bonyolult feladat, a megfigyelési terület hangszórókkal való körbezárását igényeli.
A fizikai hangtér-reprodukció alapfeladata ezek után az egyes hangszórók vezérlőjelének, vagy \emph{vezérlő függvényének} meghatározása úgy, hogy a teljes hangszórósokaság eredő hangtere megegyezzen az előírt \emph{virtuális térrel}.
A hangszórósokaság elnevezése a továbbiakban a \emph{másodlagos forráseloszlás (secondary source distribution)}.
A hangtér-reprodukció általános elrendezése \azr{fig:introduction:sfs_aim}. ábrán látható.

\begin{figure}  
\small
  \begin{minipage}[c]{0.64\textwidth}
	\begin{overpic}[width = 1\columnwidth ]{figs/sfs_aim.png}
	\small
	\put(58,30){megfigyelési terület}
	\put(0,45){\parbox{.5in}{virtuális forrás}}
	\put(45,7){másodlagos forráseloszlás}
	\end{overpic}   \end{minipage}\hfill
	\begin{minipage}[c]{0.3\textwidth}
    \caption{Általános hangtérszintézis elrendezés: a szintézis célja egy tetszőleges virtuális forráseloszlás terének előállítása a megfigyelési terület határán elhelyzetett hangszórósokaság megfelelő vezérlésével.}
\label{fig:introduction:sfs_aim}  \end{minipage}
\end{figure}

Az elmúlt néhány évtized során számos hangtér-reprodukciós technika született mind numerikus, mind analitikus megoldásokat kínálva a problémára.
Az analitikus módszerek két fő csoportra oszthatók: a szintézisproblémát leíró integrálegyenletek közvetlen megoldására szolgáló \emph{explicit}, illetve a Huygens elven alapuló ún. \emph{implicit} módszerekre.
Az utóbbi, implicit módszer, elterjedt nevén a \emph{hangtérszintézis (Wave Field Synthesis, WFS)} a jelen disszertáció fő témája.

\subsection{A hangtérszintézis története}
%
A hangtérszintézis elméleti alapjai a Delfti Műszaki Egyetemről származnak, ahol elsőként ültették át a szeizmikus hullámtan alapkoncepcióit a hangszóró- és mikrofonsorok elméletébe.
A bevezetett \emph{tradicionális hangtérszintézis} elméletének alapja a Huygens-elv matematikai formájának tekinthető Rayleigh-integrál, amely integrál egy tetszőleges hangteret egy felületi integrál formájában ad meg.
Az integrál magában tartalmazza egy sík mentén elhelyezett hangszórósokaság vezérlőfüggvényét.
A gyakorlatban is megvalósítható vonalmenti hangszórósorok vezérlőjelének meghatározásához a Rayleigh-integrált az ún. \emph{stacionárius fázis módszerével} egy kontúrintegrállá redukálták, amely implicite tartalmazza a keresett vezérlőfüggvényeket.

A technika korai megfogalmazása dipól tulajdonságokkal rendelkező hangszórósorokhoz szolgáltatott vezérlőjeleket,\footfullcite{Berkhout1988}\textsuperscript{,}\footfullcite{Berkhout1993:Acoustic_control_by_WFS} amelyet rövidesen a dinamikus hangszórókat jobban modellező másodlagos monopólusokra is kiterjesztettek.\footfullcite{doi:10.1121/1.404755}\textsuperscript{,}\footfullcite{Start1997:phd}\textsuperscript{,}\footfullcite{Verheijen1997:phd}
A technikával lehetővé vált virtuális 3D pontforrások hullámfrontjának fizikai visszaállítása a hangszórósort és a megfigyelőt tartalmazó horizontális síkban, a \emph{szintézis síkjában}.
Jól ismert volt az a tény is, hogy az ideálisan sík hangszóróeloszlás hangszórósorrá való redukálása azt eredményezi, hogy az amplitúdóhelyes szintézis még a szintézis síkjában is csak egy megadott görbe mentén, a \emph{referenciagörbén} valósítható meg.
A tradicionális elméletben ezt a görbét a hangszórósorral párhuzamos \emph{referenciavonalnak} választották meg.\footfullcite{Start1997:phd}

Már a korai hangtérszintézis elmélet is számos K+F projektnek szolgált alapul. 
Ezek közül a legkiemelkedőbb a nemzetközi CARROUSO projekt volt, amelynek célja a hangtérszintézis elméletének az abban az időben bevezetett MPEG-4 szabványba való integrálása volt.
Ez a törekvés végül nem valósult meg, azonban a projektnek több spin-off cége, így az IOSONO (ma már Barco) és a Sonic Emotion, jelenleg is gyárt és telepít egyedi tervezésű hangtérszintézis rendszereket.

A hangtérszintézis elméletének legutóbbi mérföldköve Spors és Ahrens munkájához köthető, akik a klasszikus WFS elméletet újragondolva általánosították azt tetszőleges virtuális tér és hangszórókontúrok felé.\footfullcite{Spors2008:WFSrevisited}
Az így kapott vezérlőfüggvényekkel egy tetszőleges, de vertikális irányban invariáns 2D hangtér visszaállíthatóvá vált az amplitúdóhelyes szintézist a megfigyelői tér egy pontjára, a \emph{referenciapontra} korlátozva.
Ez az általánosítás már lehetővé tette összetettebb hangterek szintézisét is, így pl. egy mozgó forrásét,\footfullcite{Ahrens2008moving} azonban a módszer  3D virtuális tér reprodukciója esetén még a referenciapontban is amplitúdóhibákhoz vezetett.
Ezen felül a két bemutatott módszer pontos kapcsolata sem volt a korábbiakban tisztázva.

\subsection{A kutatás célja}

A jelen disszertáció a kiindulási alapoktól felülvizsgálja a hangtérszintézis elméletét. 
A kutatás célja az elmélet teljes általánosítása volt, amely lehetővé teszi egy tetszőleges virtuális tér reprodukcióját tetszőleges hangszórókontúr alkalmazásával és az amplitúdóhelyes szintézis szabadon választott referenciagörbére való optimalizálásával.
Az így kapott elméleti keretrendszer speciális esetekként magában tartalmazza a korábbiakban bemutatott módszereket: 
Egy virtuális pontforrás egyenes hangszórósorral és egyenes referenciavonallal való szintézise közvetlenül a tradicionális hangtérszintézis eredményére vezet.
Hasonlóan, egy 2D virtuális teret egy referenciapontban reprodukálva a Spors által bemutatott vezérlőfüggvényeket kapjuk. 

További kérdéseket vet fel az implicit hangtérszintézis és az explicit megoldások kapcsolata is, amelyet eddig a vonatkozó irodalomban csak speciális virtuális forrás-hangszórókontúr esetekre vizsgáltak:
Kör és egyenes másodlagos forráseloszlás geometriák esetén bizonyos virtuális terekre megmutatható, hogy a hangtérszintéis az explicit, direkt megoldás nagyfrekvenciás közelítésének tekinthető.\footfullcite{Ahrens2012}\textsuperscript{,}\footfullcite{Spors2010:analysis_and_improvement}
A kutatás további célja ennek a kapcsolatnak az általánosítása volt tetszőleges virtuális térre és hangszóró-, illetve referenciakontúrokra.

A bemutatott eredmények egy összetett alkalmazási példájaként mozgó források szintézise is bemutatásra kerül.
A mozgó hangforrások által keltett hullámok reprodukciójának kérdése már szinte a hangtérszintézis bevezetésével egyidőben természetes igényként jelent meg.
A korai megoldások a tradicionális vezérlőfüggvények alapján próbálták megoldani a feladatot, a virtuális forrás pozíciójának idővariáns változtatásával.
Ez az egyszerű megoldás azonban nem képes a Doppler-hatást megfelelően visszaállítani, és így hallható torzulásokhoz vezet a szintetizált térben.
Ahrens az újragondolt hangtérszintézis elméletet terjesztette ki mozgó forrásokra is.\footfullcite{Ahrens2008moving}\textsuperscript{,}\footfullcite{Ahrens2008moving_b}
A módszer alapvetően működőképes volt, azonban a technika alapvető korlátai miatt a visszaállított tér amplitúdóját a referenciapontban sem lehetett kontrollálni.
Épp ezért, a kutatás további célja a bemutatott elméleti keretrendszer kiterjesztése volt dinamikus, idővariáns virtuális terekre.

\section{Módszer}

A következő fejezetben a disszertációban gyakran alkalmazott módszerek és analíziseszközök kerülnek rövid bemutatásra.

Az általánosított hangtérszintézis elmélet, valamint az explicit megoldás tértartománybeli alakja (amely a későbbiekben összehasonlíthatóvá teszi az explicit módszert a hangtérszintézissel) a különböző perem- és spektrális integrálok aszimptotikus közelítésén alapszik.
Munkám során bevezettem a hangterek lokális tulajdonságait leíró mennyiségeket, így a lokális hullámszámvektort és a hullámfrontgörbületet, amelyek a későbbiekben nagyban elősegítették a fenti integrálközelítések fizikai interpretációját és megértését.
Ezen lokális tulajdonságok ismert mennyiségek a nagyfrekvenciás akusztika, illetve sugárakusztika területein, de ez idáig nem kerültek bevezetésre a hangtér-reprodukció területén.

A disszertáció központi matematikai módszere az ún. \emph{stacionárius fázis módszere (stationary phase approximation, SPA)}.
Az eljárás komplex értékű függvények integráljának közelítő kiértékelésére szolgál az integrálás útjában található krtikus pontok, az ún. \emph{stacionárius pontok} körül.
Gyakorlatilag a disszertáció ezen közelítés alkalmazásáról szól, azt 
\begin{itemize}
\item a virtuális teret leíró peremintegrálokra alkalmazva, amely segítségével a peremintegrálból a vezérlőfüggvények kiolvashatók,
\item különböző hangterek és a vezérlőfüggvények Fourier integráljára alkalmazva, amely alapján meghatározhatóak az itt bevezetett tértartománybeli vezérlőfüggvények, valamint az itt bemutatott átlapolódásgátló szűrési módszer.
\end{itemize}

Az analitikus eredmények validálása a szintetizált tér numerikus szimulációjával történik, mely az egyes hangszóróelemek analitikusan rendelkezésre álló (a Green-függvény alapján) tereinek a vezérlőfüggvénnyel súlyozott összegeként állítható elő. 

\section{Eredmények}

A következő fejezet a disszertáció legfontosabb eredményeit foglalja röviden össze.
Az eredmények összefoglalásáról nemzetközi kollokvium előadás és cikk is született \cite{Firtha2019:daga_booklet}.
%Az új tudományos eredmények a következő fejezetben kerülnek részletes felsorolásra.

\subsection{Az általánosított hangtérszintézis elmélet}

A hangtér-repodukció alapelrendezése az \ref{fig:introduction:sfs_aim}. ábrán látható.
Matematikailag a szintetizált hangteret a hangszóró-vezérlőfüggvény és az egyes hangszóróelemek terét leíró Green-függvény a másodlagos forráseloszlás kontúrja fölött vett konvolúciójaként írhatjuk fel.
A vezérlőfüggvények meghatározása tehát egy dekonvolúciós, inverz probléma.

A hangtérszintézis célja a virtuális, reprodukálandó hangtér egy olyan peremintegrál formájában történő előállítása, amely peremintegrál implicite tartalmazza a keresett vezérlőfüggvényeket.
Mivel általánosan egy 3D hangtér egy zárt felületi integrál formájában adható meg, az alapvető feladat ezen felületi integrál kontúrintegrállá redukálása.
Ez a horizontális síkba való redukálás egy vertikális irányú stacionárius fázisú közelítéssel valósítható meg.

A tradicionális elmélet kiinduló pontja a Rayleigh-integrál volt, amely azonban a másodlagos forráseloszlást eleve végtelen egyenes vonalra korlátozta.
Nagyfrekvenciás feltételezéssel élve megfelelő választás a virtuális tér Kirchhoff-közelítéssel való reprezentációja, amely egy tetszőleges 3D hangteret egy a Green-függvénnyel vett konvolúciós felületi integrálként ír le.
A Kirchhoff-közelítés stacionárius fázisú kiértékelésével a felületi integrált kontúrintegrállá redukáltam, amely már magában tartalmazza a keresett vezérlőfüggvényeket.
%
\begin{figure}[t!]
\small
  \begin{minipage}[c]{0.64\textwidth}
	\begin{overpic}[width = 1\columnwidth ]{figs/25D_WFS_general.png}
	\small
	\put(2,53){(a)}
	\put(2,1){(b)}
	\end{overpic}   \end{minipage}\hfill
	\begin{minipage}[c]{0.35\textwidth}
    \caption{3D pontforrás szintézise állandósult állapotban. A virtuális forrás pozíciója $\vxs = \posvec{3}{0.4}{2.5}{0}$, gerjesztőjele harmonikus jel $f_0 = 1.5~\mathrm{kHz}$ frekvenciával.
	Az ábra (a) része a szintetizált tér valós részét, (b) része a célhangtértől vett amplitúdóhibát ábrázolja logaritmikus skálán.
	Az aktív másodlagos forrásokat folytonos, az inaktív forrásokat szaggatott fekete vonal jelzi.
	Az egyes másodlagos források referenciapozíciója a referenciagörbén jelen példában numerikusan van számítva.
%	Obviously, in the present geometry there exist secondary sources for which no unique reference position can be found.
	}
\label{fig:SFS_theory:25D_WFS_generals}   \end{minipage}
\end{figure}  

Ez a dimenziócsökkentési eljárás azonban az amplitúdóhelyes szintézis helyét egyetlen megfigyelési pontra, az ún. \emph{referenciapontra} korlátozza.
A Spors-féle hangtérszintézis elmélet ebben a referenciapontban tette lehetővé 2D virtuális terek szintézisét.
Disszertációmban bemutatom, hogy egy további, a horizontális síkban történő aszimptotikus közelítés alkalmazásával minden egyes másodlagos forráselemhez (hangszóróhoz) saját referenciapont rendelhető.
Ennek a forráselemenkénti referenciapontnak a helye analitikusan meghatározható az újonnan bevetett lokális hullámszámvektor segítségével, amely vektor definíciószerűen adott pontban a hullám lokális terjedési irányába mutat.
Rámutattam, hogy az egyes forráselemek referenciapontjának a forráselemtől vett irányát a virtuális tér lokális hullámszámvektora határozza meg, míg a forráselemtől vett távolság egy frekvenciafüggetlen erősítési tényező segítségével tetszőlegesen befolyásolható.
Teoretikus, folytonos másodlagos forráseloszlást feltételezve a referenciapontok összessége is folytonos görbét, ún. \emph{referenciagörbét} alkot.
Amennyiben a kívánt referenciagörbe alakja analitikusan ismert, meghatározható az a frekvenciafüggetlen erősítési tényező, ami az amplitúdóhelyes szintézist erre a görbére valósítja meg.
Ez alapján tetszőleges virtuális forrás-másodlagos forráseloszlás-referenciagörbe specifikus vezérlőfüggvények vezethetők be.\cite{Firtha2016_booklet}

Egyszerű példaként \azr{fig:SFS_theory:25D_WFS_generals}.\ ábra egy 3D pontforrás általános alakú másodlagos forráseloszlással történő szintézisét mutatja be.
A referenciagörbe egy koncentrikus kontúr a megfigyelési területen belül.
Az ábra (b) része alátámasztja, hogy a vezérlőfüggvényekkel a teljes megfigyelési területen fázishelyes szintézis érhető el (azaz a hullámfront alakja helyesen visszaállítható), míg az amplitúdóhelyes szintézis a referenciagörbe azon részére korlátozódik, amelyre rálátni az adott virtuális forráspozícióból.

\subsection{A tértartománybeli explicit megoldás}

Az \emph{explicit megoldás} célja a hangtér-reprodukciós inverz probléma közvetlen, direkt megoldása a konvolúciós integrál spektrális tartományba való transzformációjával.
A spektrális vagy hullámszámtartományban a konvolúció szorzattá transzformálódik és így a vezérlőfüggvény spektruma a virtuális hangtér és a másodlagos források terének a referenciagörbén mért spektrumainak hányadosaként áll elő.
A másodlagos forráseloszlást és a referenciagörbét és végtelen egyeneseknek választva a módszerhez szükséges transzformáció egy egyszerű térbeli Fourier transzformáció, és a technikát a \emph{spektrális osztás módszerének (Spectral Division Method, SDM)} nevezzük.
Természetesen a vezérlőfüggvények közvetlenül alkalmazható tértartománybeli alakja egy térbeli inverz Fourier transzformáció formájában adott, amely transzformáció azonban általában nem végezhető el analitikusan.

A stacionárius fázis módszer alkalmazásával aszimptotikus, nagyfrekvenciás közelítést adtam az explicit vezérlőfüggvény térbeli alakjára tetszőleges virtuális hangtér szintéziséhez \cite{Firtha2017:daga_booklet}.
A Kirchhoff-közelítés érvényessége mellett ezek a vezérlőfüggvények tetszőleges konvex hangszórósokaságra érvényesek (azt lokálisan lineárisnak feltételezve).
Fontos különbség a hangtérszintézishez képest, hogy az új explicit vezérlőfüggvény kiértékeléséhez a virtuális hangtér a referenciagörbén vett értékének az ismerete szükséges, míg a hangtérszintézis vezérlőfüggvények a virtuális tér a hangszórósoron mért jellemzői alapján számíthatók.

A Rayleigh-integrál aszimptotikus kiértékelésének segítségével kifejeztem az új explicit vezérlőfüggvényeket a virtuális tér a másodlagos forráseloszláson vett értéke alapján, amely eredmény egybeesik a hangtérszintézis vezérlőfüggvényekkel.
Így tehát általánosan bizonyítottam, hogy a hangtérszintézis az explicit megoldás aszimptotikus, nagyfrekvenciás közelítését adja tetszőleges virtuális forrás, másodlagos forráseloszlás és referenciagörbe esetén \cite{Firtha2018:WFS_vs_SDM_booklet}.
A két módszer aszimptotikus egyezése a vonatkozó irodalomban eddig csak speciális virtuális forrásmodellek és hangszóróeloszlások esetére volt belátva \cite{Schultz2019:HOA_vs_WFS_booklet}.

\begin{figure}[t!]  
\small
  \begin{minipage}[c]{0.6\textwidth}
	\begin{overpic}[width = 1\columnwidth ]{figs/Antialiased_synth.png}
	\small
	\put(2,53){(a)}
	\put(2,1){(b)}
	\end{overpic}   \end{minipage}\hfill
	\begin{minipage}[c]{0.35\textwidth}
    \caption{
	Sávkorlátozott impulzus gerjesztésű 3D virtuális pontforrás szintézise diszkrét másodlagos forráseloszlással, a másodlagos források közötti távolságot $\Delta x = 10~\mathrm{cm}$-re választva.
	Az (a) ábra a forráseloszlás diszkretizálásának hatását mutatja be, ami átlapolódó visszhangok megjelenéséhez vezet a virtuális hullámfront mögött.
	A (b) ábra az átlapolódásgátló szűrés hatását mutatja.
	A szűrés eredményeképp egy kitüntetett irányban (az ábrán nyíllal jelölve) az átlapolódó hullámfrontok teljesen megszüntethetőek.
   % The arrow originates at the SSD element with no angular bandwidth limitation, i.e. performing full-band synthesis.
   }
\label{fig:SFS_theory:anti-aliased_synthesis}   \end{minipage}
\end{figure} 

\vspace{3mm}
A hangtérszintézis gyakorlati alkalmazása során az ideálisan folytonos másodlagos pontforráseloszlás helyett egymástól véges távolságra, diszkrét pozíciókban elhelyezett hangszórósokaságot alkalmazunk.
Az idealizált fizikai modell megsértéséből fakadóan a szintetizált térben járulékos hatások, ún. \emph{térbeli átlapolódási jelenségek} jelennek meg.
A jelenség az időtartományban az eredeti, virtuális hullámfrontot követő másodlagos, az egyes hangszóróelemekből származó gömbi hullámfrontok sorozataként jelentkezik.
Ezek az átlapolódó ,,visszhangok'', amelyek az egyes hangszóró vezérlőfüggvények az \emph{átlapolódási frekvencia} feletti komponenseit tartalmazzák, \azr{fig:SFS_theory:anti-aliased_synthesis}. (a) ábrán figyelhetők meg.

Az explicit hangtér-reprodukciós megoldás nagy előnye, hogy a fenti átlapolódási jelenségek a spektrális tartományban egyszerűen modellezhetőek, amely modell a módszerek ekvivalenciája miatt alkalmazható a hangtérszintézisre is.
Az explicit vezérlőfüggvények aszimptotikus vizsgálatával azonosítottam, mely másodlagos forráselemek járulnak hozzá a térbeli átlapolódáshoz egy adott frekvencián.
Adott frekvencián ezeket a forrásokat elnémítva a térbeli átlapolódás elkerülhető, a megfigyelési terület csökkenése árán.
Szélessávú gerjesztés esetén a fenti átlapolódásgátló szűrés a vezérlőfüggvények pozíciófüggő aluláteresztő szűrésével valósítható meg.
Végeredményben tehát egyszerű aluláteresztő szűrő alapú átlapolódásgátló szűrési stratégiát vezettem be, ahol a szűrő vágási frekvenciája a virtuális hangtér lokális hullámszámvektorán keresztül analitikusan adott.
A módszer alkalmazásával a megfigyelési terület egy kitüntetett iránya mentén átlapolódásmentes, széles sávú szintézis valósítható meg, ahogy az \azr{fig:SFS_theory:anti-aliased_synthesis}. (b) ábra mutatja.
A tér többi irányában a virtuális hullámfront sávkorlátozott a szűrés miatt, valamint oldalirányba terjedő átlapolódó hullámfrontok maradnak jelen a térben, amelyek azonban már kizárólag irányított másodlagos forrásokkal nyomhatóak el.

A bevezetett módszer alternatívájaként a probléma megközelíthető pusztán a tértartományban, szintén a lokális hullámszámvektor alkalmazásával.
Ezt a megközelítést alkalmazza a \cite{8611109_booklet}-ban leírt módszer, amely folyóiratcikknek társszerzője vagyok.

\subsection{Mozgó források szintézise}

A bevezetett elméleti keretrendszer alkalmazására szolgáló összetett példaként megvizsgáltam a mozgó hangforrások által keltett tér reprodukciójának lehetőségeit.
Ebben az esetben az elsődleges feladat a \emph{Doppler-hatás} pontos visszaállítása, amely azonban eleve biztosítva van megfelelő analitikus forrásmodell alkalmazása esetén.

Munkám során kiterjesztettem a lokális hullámszámvektor fogalmát idővariáns hangterekre is és a stacionárius fázisú módszerrel mozgó források szintézisére alkalmas hangtérszintézis vezérlőfüggvényt vezettem be.
A vezérlőfüggvény lehetővé teszi egy tetszőleges (a megfigyelő síkjában) mozgó forrás szintézisét tetszőleges konvex másodlagos forráseloszlás alkalmazásával. 
Az amplitúdóhelyes szintézis tetszőleges referenciagörbén megvalósítható \cite{firtha2016wave_booklet, doi:10.1121/1.4996126_booklet, Firtha2015:daga_booklet}.
A módszer alkalmazását \azr{fig:SFS_theory:25D_WFS_moving_source}. ábra szemlélteti egy egyenletesen mozgó harmonikus forrás szintézisének példáján keresztül.
Látható, hogy a vezérlőfüggvények a reprodukált tér amplitúdóját az előírt referenciagörbére optimalizálják. 

\begin{figure}[t!]
\small
  \begin{minipage}[c]{0.64\textwidth}
	\begin{overpic}[width = 1\columnwidth ]{figs/25D_WFS_ms.png}
	\small
	\put(2,53){(a)}
	\put(2,1){(b)}
	\end{overpic}   \end{minipage}\hfill
	\begin{minipage}[c]{0.35\textwidth}
    \caption{
	Egyenletesen mozgó harmonikus, $f_0 = 1.5~\mathrm{kHz}$ frekvenciájú pontforrás szintézise általános hangszórógörbe alkalmazásával.
	A forrás sebessége $|\mathbf{v}| = 150~\mathrm{m/s}$.
	A referenciagörbe egy koncentrikus kontúr a megfigyelési területen.
	Az (a) ábra a szintetizált tér valós részét, a (b) ábra a szintézis abszolút hibáját mutatja be.
	A másodlagos forráseloszlás aktív részét a pillanatkép rögzítésének pillanatában folytonos, az inaktív részét szaggatott fekete vonal jelzi.
    }
\label{fig:SFS_theory:25D_WFS_moving_source}  \end{minipage}
\end{figure}

Az egyenes pályán egyenletesen mozgó források speciális esetében a forrás által lesugárzott tér spektrális, hullámszámtartomány-beli leírása analitikusan is a rendelkezésre áll.
Ez a leírás lehetővé tette explicit vezérlőfüggvények meghatározását mozgó források szintéziséhez \cite{Firtha2014:daga_booklet, Firtha2014:isma_booklet}.
Az így bevezetett vezérlőfüggvény referenciamegoldást szolgáltat a problémára, mivel az alkalmazásával közelítések nélküli, tökéletes reprodukció érhető el a referenciavonal mentén.
Ezen felül a hullámszámtartomány-beli leírás lehetővé teszi a térbeli átlapolódási jelenségek tárgyalását.
Dolgozatomban megmutatom, hogy a mozgó források szintézisét megvalósító hangtérszintézis vezérlőfüggvények---hasonlóan a statikus esethez---nagyfrekvenciás közelítését adják az explicit megoldásnak.

\begin{figure}
	\centering
	\begin{overpic}[width = 1\columnwidth]{figs/antialiased_synth_moving_source.png}
	\footnotesize
	\put(0, 0){(a)}
	\put(50,0){(b)}
	\end{overpic}
\caption{Általános, görbe pálya mentén mozgó forrás szintézise, $\Delta x = 10~\mathrm{cm}$ osztásközű, $R_{\mathrm{SSD}} = 2~\mathrm{m}$ sugarú kör alakú másodlagos forráseloszlással.
	A forrás sebessége állandó, $|\mathbf{v}| = \frac{3}{4}c$.
	Az (a) ábra a forrás diszkretizálás hatását szemlélteti, jól látható átlapolódási jelenségekkel.
	Az legnagyobb energiájú átlapolódó hullámok a mozgó forrás előtti másodlagos forrásokból származnak: itt az lokális, érzékelt frekvencia megnő a Doppler-hatás miatt.
	A (b) ábra a bevezetett átlapolódásgátló szűrés hatását mutatja be.}
	\label{Fig:Moving_sources:anti-aliased_synth_moving_source}
\end{figure}

Végül kutatásom célja a mozgó források szintézise esetén jelentkező átlapolódási jelenségek vizsgálata, valamint a bevezetett átlapolódásgátló szűrés erre az esetre való kiterjesztése volt.
Ebben az esetben az átlapolódási jelenségek a szintetizált tér fokozottan hallható torzulásához vezet, mivel az átlapolódó másodlagos hullámfrontok más Doppler-tolást szenvednek, mint az ideális, virtuális hullámfront.
Ennek eredményeképp a statikus esetben is fellépő \emph{hangszíneződésen} felül nemkívánatos frekvenciatorzítás is megjelennek a szintetizált térben \cite{firtha2016:daga_booklet}.

Harmonikus forrásgerjesztés esetén ez a frekvenciatorzítás új, nemkívánatos frekvenciakomponensekként jelenik meg, ahogy az már korábban is ismert volt a vonatkozó irodalomban.\footfullcite{Franck2007}
A szintetizált tér spektrális vizsgálata alapján analitikus formulát adtam ezekre az átlapolódó frekvenciakomponensek frekvenciájára, és rámutattam, hogy a másodlagos forráseloszlás alakjának megfelelő megválasztásával a frekvenciatorzítás hatása jelentősen csökkenthető \cite{firtha2016:daga_booklet}.
Az átlapolódás elkerülése szempontjából az optimális választás kör alakú hangszórósokaság alkalmazása.

Végezetül kiterjesztettem a korábban bevezetett átlapolódás elleni stratégiát mozgó források szintézisére is.
A javasolt módszerrel az átlapolódás a hangszóró vezérlőfüggvények egyszerű aluláteresztő szűrésével megvalósítható.
A szűrő idővariáns vágási frekvenciája analitikusan ismert, a mozgó forrás lokális hullámszámvektorának függvényében \cite{Firtha2018_daga_moving_source_booklet}.
Hasonlóan a statikus esethez, az átlapolódásmentes, teljessávú szintézis csupán egy kitüntetett irányban valósítható meg.
Kör alakú hangszóróeloszlás esetén ez mindig a radiális irány, így a hangszórótömb közepén tetszőleges időpillanatban tetszőleges forráspálya mellett optimális szintézis érhető el.
Ezt az esetet szemlélteti \azr{Fig:Moving_sources:anti-aliased_synth_moving_source}. ábra egy általános görbe pálya mentén mozgó impulzussorozat-gerjesztésű forrás szintézise esetén.

\clearpage	
\section{Tézisek}
\label{sec:theses}

\begin{thesisgroup}[A hangtérszintézis elméletének általánosítása]
Bevezettem egy általánosított hangtérszintézis elméletet, amely lehetővé teszi egy tetszőleges virtuális tér reprodukcióját tetszőleges alakú, konvex hangszóróeloszlás alkalmazásával és a szintézis amplitúdója tetszőleges alakú referenciagörbére optimalizálható.
Az általánosított elmélet speciális esetekként tartalmazza  a korábbi hangtérszintézis módszereket \cite{Firtha2016_booklet}.
\begin{thesis}
Szemléletes fizikai interpretációt adtam a peremintegrálok stacionárius fázisú közelítésének:
Bevezettem a tetszőleges harmonikus hangtérre értelmezett lokális hullámszámvektort, amely felhasználásával megmutattam, hogy a peremintegrál stacionárius fázis közelítése a megfigyelési pozícióban a virtuális hangtér és a perem mentén elhelyezett másodlagos források terének hullámfront-illesztéseként értelmezhető.
\end{thesis}
\begin{thesis}
Hangtérszintézis vezérlőfüggvényeket vezettem be tetszőleges konvex másodlagos forráseloszlásra a fenti fizikai magyarázat alapján, nagyfrekvenciás feltételezések mellett.
\end{thesis}
\begin{thesis}
Analitikus összefüggést adtam a referenciagörbe általános alakjára, amely görbe azon pontok összessége, amelyeken az szintézis amplitúdóhibája minimális.
Bemutattam, hogyan befolyásolható ezen referenciagörbe alakja a vezérlőfüggvény frekvenciafüggetlen amplitúdótényezőjével.
Bemutattam, hogy a korábbi hangtérszintézis módszerek milyen speciális esetei az általánosított elméletnek virtuális forrás-másodlagos eloszlás-referenciagörbe szempontjából.
\end{thesis}
\end{thesisgroup}

\begin{thesisgroup}[Térbeli explicit megoldás és ekvivalencia a hangtérszintézissel]
A hangtérszintézis mellett---amely célja egy olyan kontúrintegrál megfogalmazása, amely a keresett hangszóró vezérlőfüggvényeket implicite tartalmazza---egyszerű geometriák esetén explicit vezérlőfüggvények számítása is lehetséges a probléma térbeli frekvenciatartományba transzformációja alapján.
Egyenes hangszóróeloszlás esetén az explicit vezérlőfüggvény a virtuális tér és a másodlagos forráselemek fix referenciavonalon mért spektrális hányadosának térbeli inverz Fourier transzformáltjaként áll elő.
Az implicit és explicit módszerek közötti kapcsolatot a korábbiakban csak speciális virtuális hangterekre vizsgálták.
Az explicit megoldás aszimptotikus közelítésével nagyfrekvenciás térbeli explicit vezérlőfüggvényeket vezettem be, és rámutattam az implicit és explicit megoldások általános nagyfrekvenciás ekvivalenciájára \cite{Firtha2017:daga_booklet, Firtha2018:WFS_vs_SDM_booklet}
\begin{thesis}
Tértartománybeli explicit vezérlőfüggvényeket vezettem be a térbeli Fourier transzformáció stacionárius fázisú közelítésével.
Az új vezérlőfüggvények kiértékeléséhez a virtuális hangtér jellemzőit az előírt referenciagörbén szükséges ismerni \cite{Firtha2017:daga_booklet}.\end{thesis}
\begin{thesis}
Bebizonyítottam, hogy nagyfrekvenciás feltételek mellett az explicit megoldás és a hangtérszintézis (azaz az implicit megoldás) aszimptotikusan megegyeznek tetszőleges virtuális hangtér esetére \cite{Firtha2018:WFS_vs_SDM_booklet}.
A bizonyítás alapja a térbeli explicit vezérlőfüggvények kifejezése volt a virtuális forrás másodlagos forráseloszláson mért jellemzői alapján.
\end{thesis}
\begin{thesis}
Egyszerű átlapolódásgátló stratégiát mutattam be a másodlagos forráseloszlás diszkretizálásából származó átlapolódó hullámfrontok csillapítása érdekében.
A módszer a térbeli Fourier transzformáció aszimptotikus közelítésén alapszik.
A bemutatott stratégia a vezérlőfüggvények egyszerű aluláteresztő szűrésével implementálható \cite{Firtha2018_daga_a_booklet}.
\end{thesis}
\end{thesisgroup}

\begin{thesisgroup}[Mozgó források hangtérszintézise]
Dinamikusan változó akusztikai környezetek reprodukciója szempontjából a mozgó források terének megfelelő szintézise alapvető feladat.
Kiterjesztettem az itt bemutatott hangtérszintézis elméletet mozgó források szintézisére is.
\begin{thesis}
Kiterjesztettem a 3D hangtérszintézis elméletet tetszőleges pályán, tetszőleges sebességprofillal mozgó 3D pontforrások szintézisére.
A megoldás a Doppler-hatást analitikusan figyelembe veszi \cite{Firtha2015:daga, firtha2016wave_booklet, doi:10.1121/1.4996126_booklet}.\end{thesis}
\begin{thesis}
Hangtérszintézis vezérlőfüggvényeket határoztam meg egy tetszőleges alakú 2D hangszórókontúrra, amellyel egy tetszőleges---a hangszórók síkjában---mozgó  pontforrás tere szintetizálható, az amplitúdóhelyes szintézist tetszőleges referenciagörbére megvalósítva \cite{doi:10.1121/1.4996126_booklet}.
A megoldás alapja a stacionárius fázis módszer kiterjesztése a dinamikus, idővariáns esetre.
Rámutattam, hogy lineáris hangszórósor és ezzel párhuzamos referenciavonal esetén az általános megoldás a tradicionális pontforrás vezérlőfüggvényeket eredményezi a statikus forráspozíciót az emissziós időben felvett pozícióra cserélve \cite{doi:10.1121/1.4996126_booklet}.
\end{thesis}
\begin{thesis}
Zárt, analitikus vezérlőfüggvényt adtam egyenes vonalon egyenletesen mozgó pontforrás szintézisére, amely esetben a vezérlőfüggvényben jelen lévő emissziós idő, illetve terjedési késleltetések explicite kifejezhetőek \cite{firtha2016wave_booklet}.
\end{thesis}
\begin{thesis}
Hangtérszintézis vezérlőfüggvényeket határoztam meg egyenes hangszórósorra a frekvenciatartományban a mozgó források spektrális leírásának stacionárius fázisú közelítésével \cite{firtha2015sound_booklet}.
\end{thesis}
\end{thesisgroup}

\begin{thesisgroup}[Explicit megoldás mozgó források szintézisére]
Analitikus kifejezést adtam egyenes vonal mentén egyenletesen mozgó pontforrás terének térbeli Fourier transzformáltjára.
Mivel az explicit megoldás nincs időinvariáns terekre korlátozva, ezért a kifejezés felhasználható explicit vezérlőfüggvények számítására mozgó forrásokhoz.
\begin{thesis} 
Explicit vezérlőfüggvényeket határoztam meg egyenletesen mozgó források reprodukciójához a hullámszámtartományban.
A másodlagos forráseloszlással párhuzamosan mozogó források speciális esetében az explicit vezérlőfüggvényt analitikusan meghatároztam a tér-frekvenciatartományban is \cite{Firtha2014:daga_booklet, Firtha2014:isma_booklet, firtha2015sound_booklet}.
A frekvenciatartománybeli hangtérszintézis vezérlőfüggvények alapján rámutattam, hogy---a statikus esethez hasonlóan---az implicit megoldás az explicit megoldásnak a nagyfrekvenciás, aszimptotikus közelítését adja \cite{firtha2015sound_booklet}. 
\end{thesis}
\begin{thesis}
Analitikus leírását adtam a térbeli átlapolódási jelenségekre mozgó források szintézise esetében.
Rámutattam, hogy a nemkívánatos frekvenciakomponensek a Green-függvény spektrumában jelenlévő szinguláris pontokhoz köthetők, és ez alapján analitikus összefüggést adtam ezen nemkívánatos összetevők frekvenciájának meghatározására.
Bemutattam, hogy ezek a hatások csökkenthetők megfelelően megválasztott másodlagos forráseloszlás alkalmazásával \cite{firtha2016:daga_booklet}.
\end{thesis}
\begin{thesis} 
Kiterjesztettem a térbeli átlapolódásgátló szűrési stratégiát tetszőleges trajektórián mozgó források esetére.
A javasolt módszerrel az átlapolódó hullámfrontok egyszerű idővariáns aluláteresztő szűréssel kiküszöbölhetőek \cite{Firtha2018_daga_moving_source_booklet}.
\end{thesis}
\end{thesisgroup} 
\clearpage

\section*{Publikációs lista}
\addcontentsline{toc}{sec}{Publikációk}

\nocite{*}
\begin{refcontext}[labelprefix=J]
\printbibliography[title={Folyóiratcikkek}, keyword=J, heading=subbibliography] 
\end{refcontext}
\begin{refcontext}[labelprefix=C]
\printbibliography[title={Konferenciaközlemények}, keyword=C, heading=subbibliography] 
\end{refcontext}
\begin{refcontext}[labelprefix=O]
\printbibliography[title={Egyéb publikációk}, keyword=O, heading=subbibliography] 
\end{refcontext}

\end{document}
